\section{System}
\label{sec:system}

Each question passes through retrieval, an answer with citations, the check and at most one revision, and ends in one of three outcomes. \Cref{fig:pipeline} shows the path; this section takes each step in turn.

\begin{figure}[tbp]
  \centering
  \begin{tikzpicture}[
      font=\sffamily\small,
      node distance=10mm and 8mm,
      box/.style={draw=black!70, rounded corners=2pt, align=center, inner sep=4pt, minimum height=11mm},
      check/.style={box, draw=prussian, very thick},
      outcome/.style={align=center, inner sep=2pt, text=prussian},
      flow/.style={-{Stealth[length=2mm]}, black!70},
      label/.style={font=\sffamily\footnotesize},
    ]
    \node[box] (question) {Question};
    \node[box, right=of question] (search) {BM25 search\\8 excerpts};
    \node[box, right=of search] (answer) {Model answer\\with citations};
    \node[check, right=of answer] (check) {Check\\quote and numbers};
    \node[outcome, right=of check] (shown) {Verified answer\\shown};
    \node[box, below=of check] (revise) {One revision\\with the reason};
    \node[outcome, right=of revise] (withheld) {No verified\\answer};
    \node[outcome, below=of answer] (none) {The datasheet\\does not say};

    \draw[flow] (question) -- (search);
    \draw[flow] (search) -- (answer);
    \draw[flow] (answer) -- (check);
    \draw[flow] (check) -- node[label, above] {pass} (shown);
    \draw[flow] ([xshift=-5mm]check.south) -- node[label, left] {fail} ([xshift=-5mm]revise.north);
    \draw[flow] ([xshift=5mm]revise.north) -- node[label, right] {recheck} ([xshift=5mm]check.south);
    \draw[flow] (revise) -- node[label, above] {fail} (withheld);
    \draw[flow] (answer) -- node[label, left] {unanswerable} (none);
  \end{tikzpicture}
  \caption{The path of one question. Only an answer that passes the check is shown. A rejected answer is sent back once with the reason it failed, and a revision that also fails is withheld.}
  \label{fig:pipeline}
\end{figure}

\subsection{Pages and text}

The datasheets are downloaded from Bosch Sensortec and checked against a SHA-256 hash recorded in the repository, so every run reads the same files; the PDFs themselves are not redistributed. The text of each page is extracted with pdfium. Table rows come out in reading order, one row to a line, but the column a value sat in is lost: a row of minimum, typical and maximum values reads \enquote{\texttt{1.8 2.8 µA}} with no sign of which column was empty. Figures and graphs carry no text layer and are invisible to the assistant.

One normaliser prepares text for both search and the check, so the two never disagree about what a page says. It applies Unicode compatibility folding, maps the Unicode minus sign and every dash to a hyphen, unifies the two micro signs, removes a hyphen between two letters, and collapses runs of white space. The hyphen rule exists because pdfium marks a line-break hyphen with one code point whether the word was split (\enquote{consist-ing}) or genuinely hyphenated (\enquote{low-power}); removing every hyphen between letters on both sides of the comparison makes both spellings match.

\subsection{Retrieval}

Each page is cut into chunks of whole lines, merged up to about 300 words and never spanning two pages, and each chunk keeps its datasheet and page. Search ranks chunks with BM25~\cite{robertson2009}, with $k_1 = 1.2$ and $b = 0.75$, and passes the top eight, about 2,400 words, to the model. BM25 was chosen over embeddings because datasheet questions turn on exact terms such as register names, and because it is deterministic. When a question names a part, only that part's datasheet is searched. Without this rule a BMP390 question ranked BMA456 and BMI270 pages first: every page header repeats the part number, so as a search term it barely separates one datasheet from another. A question that shares no term with any chunk is declared unanswerable without a model call.

\subsection{Answer}

The model, \texttt{claude-sonnet-5-5}, receives the excerpts, each marked with its datasheet and page, and returns a structured reply: whether the excerpts answer the question, a short answer, and one to three citations, each naming a datasheet and page and quoting a passage of at least four words. The prompt asks for numbers written as the datasheet writes them, no unit conversions, and no number repeated from the question unless a quote contains it. A request the API refuses or truncates, a failed request, or a reply that does not match the schema is recorded as an error, never as a sign that the datasheet does not say.

\subsection{The check}

The check is a pure function of the answer and the page texts. An answer passes only if it cites at least one passage; every citation names a datasheet and page that exist; every quote is at least four words long and appears, after normalisation, on the page it cites; and every number in the answer appears in at least one of its quotes. The four-word minimum stops a quote such as \enquote{\qty{3.6}{\volt}}, which appears on many pages, from passing as evidence. The numbers are pooled across citations because one answer may combine values from up to three passages.

Numbers are extracted as signed decimals, scientific notation and hexadecimal (\texttt{0x76}). A digit attached to letters, as in \texttt{BME280} or \texttt{I2C}, is not a number, and an \texttt{x} between digits separates two values, so \texttt{2.5x3.0} yields 2.5 and 3.0. Decimals compare by value with their sign, so 3.30 matches 3.3 but $-40$ does not match 40; hexadecimal compares as text, ignoring case.

\subsection{Revision and outcomes}

A rejected answer is sent back to the model in the same conversation, followed by the reason it failed, for example \enquote{answer states 0, 2, which no quote contains}. The model is asked to fix the cause, either by quoting a passage that holds the number or by leaving out a number the question did not ask for, and is told not to drop the value the question asks for. One revision is allowed. The revision faces the same check, so it can move an answer from withheld to shown only by passing.

Each question ends in one of three outcomes. A verified answer is shown with its citations. \enquote{No verified answer} means every draft failed the check: the datasheet may still state the value, and the reason is shown. \enquote{The datasheet does not say} means the model found no passage that answers. A web page shows each verified answer beside the rendered datasheet page with the quoted passage highlighted, and marks a number's place on the page when the reader points at it in the answer.
